\subsection{乘法算子}

设 X 是局部紧空间，\(\mu\) 是 X 上的正测度。我们考察 \(\mathrm{L}^{2}(\mathrm{X},\mu)\) 上的乘法算子；它们是稠定的闭算子（IV，第 232 页，命题 5）。对 X 上任何 \(\mu\)-可测函数 g，我们都用 \(m_{g}\) 表示 \(\mathrm{L}^{2}(\mathrm{X},\mu)\) 上以 g 为乘子的乘法部分算子。

命题 22. — 设 \(g\) 是 X 上的 \(\mu\)-可测函数。

\begin{enumerate}
\def\labelenumi{\alph{enumi})}
\item
  \(m_{g}\) 的谱是相对于 \(\mu\) 的 \(g\) 的本质像 S；
\item
  设 \(\lambda \in \mathbf{C} - \mathrm{Sp}(m_g)\)。预解式 \(\mathrm{R}(m_g, \lambda)\) 是乘法算子 \(m_h\)，其中 \(h\) 是 X 上如下定义的函数：\(h(x) = 0\)（若 \(g(x) = \lambda\)），否则 \(h(x) = (\lambda - g(x))^{-1}\)。
\end{enumerate}

我们来证明 \(\mathbf{C} - \mathbf{S}\) 包含于 \(m_g\) 的预解集。设 \(\lambda \in \mathbf{C} - \mathbf{S}\)。存在 \(\lambda\) 的开邻域 U，使 X 的子集 \(\mathrm{Y} = g^{-1}(\mathrm{U})\) 是局部 \(\mu\)-零测的。X 上如下定义的函数 \(k\)：\(k(x) = (\lambda - g(x))^{-1}\)（若 \(x \notin \mathrm{Y}\)），而 \(k(x) = 0\)（若 \(x \in \mathrm{Y}\)），它属于 \(\mathcal{L}^{\infty}(\mathrm{X}, \mu)\)（IV，第 184 页，引理 5）；因此以 \(k\) 为乘子的乘子算子是 \(\mathrm{L}^{2}(\mathrm{X}, \mu)\) 的自同态。

由于 \(|gk| \leqslant 1 + |\lambda k|\)，我们有

\[
| g k f | \leqslant | f | + | \lambda k f |
\]

对 \(f \in \mathcal{L}^2(\mathrm{X},\mu)\) 成立，这蕴含 \(m_k\) 的像包含于 \(m_g\) 的定义域。反之，设 \(f \in \mathcal{L}^2(\mathrm{X},\mu)\)，其类 \(\widetilde{f}\) 属于 \(m_g\) 的定义域。在局部 \(\mu\)-零测集 Y 之外，我们有 \(f(x) = k(x)(\lambda - g(x))f(x)\)，故 \(\widetilde{f}\) 在 \(m_k\) 的像中。同一公式证明 \(\lambda\) 属于 \(m_g\) 的预解集且 \(\mathrm{R}(m_g,\lambda) = m_k\)。由于集合 Y 是局部 \(\mu\)-零测的，乘法算子 \(m_k\) 与算子 \(m_h\)（如断言 \(b\)）中所描述的）一致。

反之，我们来证明 \(\mathbf{C} - \mathrm{Sp}(m_g)\) 包含于 \(\mathbf{C} - \mathrm{S}\)。设 \(\lambda \in \mathbf{C} - \mathrm{Sp}(m_g)\)。取实数 \(\mathrm{M} > \| \mathrm{R}(m_g, \lambda) \|\)。用 Y 表示由这样的 \(x \in \mathrm{X}\) 组成的集合：它们满足 \(|\lambda - g(x)| < \mathrm{M}^{-1}\)。我们来证明 Y 是局部 \(\mu\)-零测的，这将蕴含 \(\lambda\) 不属于 S，并完成证明。

设 \(K\) 是 \(X\) 的紧子集。设 \(\varphi\) 是 \(Y \cap K\) 的特征函数；它是 \(\mathcal{L}^{2}(X,\mu)\) 中的元素，我们把它在 \(\mathrm{L}^{2}(X,\mu)\) 中的类记作 \(\widetilde{\varphi}\)。设 \(\psi\) 是 \(\mathcal{L}^{2}(X,\mu)\) 中的函数，它在 \(\mathrm{L}^{2}(X,\mu)\) 中的类是 \(\mathrm{R}(m_g,\lambda)(\widetilde{\varphi})\)。我们有 \(\mathrm{R}(m_g,\lambda)(\widetilde{\varphi}) \in \operatorname{dom}(m_g)\) 且 \((\lambda - m_g)(\mathrm{R}(m_g,\lambda)(\widetilde{\varphi})) = \widetilde{\varphi}\)，故 \((\lambda - g(x))\psi(x) = 1\) 对 \(\mu\)-几乎所有的 \(x \in Y \cap K\) 成立。由此可得

\[
\| \mathrm{R} (m _ {g}, \lambda) \| ^ {2} \| \widetilde {\varphi} \| ^ {2} \geqslant \| \mathrm{R} (m _ {g}, \lambda) (\widetilde {\varphi}) \| ^ {2} \geqslant \mathrm{M} ^ {2} \mu (\mathrm{Y} \cap \mathrm{K}) = \mathrm{M} ^ {2} \| \widetilde {\varphi} \| ^ {2}.
\]

鉴于 M 的取法，这表明 \(\varphi\) \(\mu\)-几乎处处为零。于是集合 Y \(\cap\) K 是 \(\mu\)-零测的。因此集合 Y 是局部 \(\mu\)-零测的（INT，IV，第 172 页，§ 5，第 2 目，命题 5）。

命题 23. --- 设 g 是 X 上的 μ-可测函数。乘法算子 \(m_{g}\) 的伴随算子是 \(m_{\overline{g}}\)。

对每个整数 \(n \geqslant 1\)，设 \(\varphi_{n}\) 是由这样的元素 \(x \in X\) 组成的集合的特征函数：它们满足 \(|g(x)| \leqslant n\)，并设 \(\widetilde{\varphi}_{n}\) 是它在 \(\mathrm{L}^{2}(\mathrm{X}, \mu)\) 中的类。设 \(f \in \mathcal{L}^{2}(\mathrm{X}, \mu)\)，其类 \(\widetilde{f}\) 属于 \(\operatorname{dom}(m_{g}^{*})\)，又设 \(\psi\) 是以 \(m_{g}^{*}(\widetilde{f})\) 为类的函数。

对每个 \(h \in \mathcal{L}^2(\mathrm{X}, \mu)\)，若其类 \(\widetilde{h}\) 属于 \(\operatorname{dom}(m_g)\)，则也有 \(\widetilde{h}\widetilde{\varphi}_n \in \operatorname{dom}(m_g)\)，从而 \(\langle \widetilde{f} | m_g(\widetilde{h}\widetilde{\varphi}_n) \rangle = \langle m_g^*(\widetilde{f}) | \widetilde{h}\widetilde{\varphi}_n \rangle\)。由此得到等式

\[
\int_ {\mathrm{X}} \overline {{(f \overline {{g}} - \psi)}} \varphi_ {n} h d \mu = 0.
\]

由于 \(m_{g}\) 的定义域在 \(\mathrm{L}^{2}(\mathrm{X},\mu)\) 中稠密，我们推出 \((f\overline{g}-\psi)\varphi_{n}\) \(\mu\)-几乎处处为零。由于 n 是任意的，这就意味着 \(m_{g}^{*}(\widetilde{f})\) 在 \(\mathrm{L}^{2}(\mathrm{X},\mu)\) 中等于 \(f\overline{g}\) 的类。特别地，由 \(m_{g}^{*}(\widetilde{f}) \in \mathrm{L}^{2}(\mathrm{X},\mu)\)，我们断定 f 属于 \(m_{\overline{g}}\) 的定义域，且 \(m_{g}^{*}(\widetilde{f}) = m_{\overline{g}}(\widetilde{f})\)。

因此 \(m_{g}\) 的伴随算子是 \(m_{\overline{g}}\) 的扩张。此外，我们有

\[
\langle f \mid m _ {g} (h) \rangle = \int_ {\mathrm{X}} \overline {{f}} \cdot (g h) d \mu = \int_ {\mathrm{X}} f \overline {{g}} h d \mu
\]

对一切 \(f \in \mathrm{L}^{2}(\mathrm{X}, \mu)\) 与 \(h \in \mathrm{dom}(m_{g})\) 成立，这就证明线性形式 \(h \mapsto \langle f | m_{g}(h) \rangle\) 当 \(m_{\overline{g}}(f) g\) 属于 \(\mathrm{L}^{2}(\mathrm{X}, \mu)\) 时是连续的。于是，\(m_{\overline{g}}\) 的定义域包含于 \(m_{g}^{*}\) 的定义域，证明至此完成。

推论。--- 设 g 是 X 上的 μ-可测函数。乘法算子 \(m_{g}\)（\(\mathrm{L}^{2}(\mathrm{X},\mu)\) 上的）是自伴的（相应地，正的），当且仅当函数 g 局部 μ-几乎处处取实值（相应地，局部 μ-几乎处处为正）。

第一个断言由该命题得出。若 \(g\) 局部 \(\mu\)-几乎处处为正，则有 \(\langle f|m_{g}(f)\rangle=\int_{\mathrm{X}}g|f|^{2}d\mu\geqslant0\)，对一切 \(f\in\mathrm{L}^{2}(\mathrm{X},\mu)\) 成立，故部分算子 \(m_{g}\) 是正的。

反之，若 \(m_{g}\) 是正的，则其谱包含于 \(\mathbf{R}_{+}\)（IV，第 248 页，命题 17）；由于该谱是 g 的 \(\mu\)-本质像（命题 22，a）），这就意味着 g 局部 \(\mu\)-几乎处处为正。

引理 6. --- 设 \(g_{1}\) 与 \(g_{2}\) 是 X 上的 \(\mu\)-可测函数。

\begin{enumerate}
\def\labelenumi{\alph{enumi})}
\item
  部分算子 \(m_{g_{1}} + m_{g_{2}}\) 是可闭的，其闭包是乘法算子 \(m_{g_{1}+g_{2}}\)；
\item
  有 \(m_{g_{1}} \circ m_{g_{2}} \subset m_{g_{1}g_{2}}\)；
\item
  设 \(g_{1}\) 有界。这时我们有 \(m_{g_{2}} \circ m_{g_{1}} = m_{g_{1}g_{2}}\)。此外，\(m_{g_{2}}\) 的定义域包含于 \(m_{g_{1}g_{2}}\) 的定义域，且 \(m_{g_{1}} \circ m_{g_{2}}\) 是 \(m_{g_{1}g_{2}}\) 到 \(\text{dom}(m_{g_{2}})\) 上的限制。
\end{enumerate}

\(m_{g_{1}+g_{2}}\) 是 \(m_{g_{1}} + m_{g_{2}}\) 的扩张，这是初等的；于是后一算子是可闭的，且 \(\overline{m_{g_{1}} + m_{g_{2}}} \subset m_{g_{1}+g_{2}}\)。

设 \(f \in \mathcal{L}^2(\mathrm{X},\mu)\)，使函数 \(h = (g_1 + g_2)f\) 属于 \(\mathcal{L}^2(\mathrm{X},\mu)\)。对每个整数 \(n \geqslant 1\)，用 \(\mathrm{X}_n\) 表示由这样的 \(x \in \mathrm{X}\) 组成的集合：它们满足 \(|g_1(x)| + |g_2(x)| \leqslant n\)，并设 \(\varphi_n\) 是 \(\mathrm{X}_n\) 的特征函数。我们有 \(\varphi_n f \in \mathrm{dom}(m_{g_1} + m_{g_2})\)。由于 \((\varphi_n f)(x)\) 收敛于 \(f(x)\)，对一切 \(x \in \mathrm{X}\)，且对每个 \(n \in \mathbf{N}\) 有 \(|\varphi_n f| \leqslant |f|\) 与 \(|(g_1 + g_2)\varphi_n f| \leqslant |(g_1 + g_2)f| = |h|\)（由假设 \(h \in \mathcal{L}^2(\mathrm{X},\mu)\)），勒贝格定理（INT，IV，第 137 页，§ 3，第 7 目，定理 6）蕴含：\(\mathrm{L}^{2}(\mathrm{X},\mu)\times\mathrm{L}^{2}(\mathrm{X},\mu)\) 中 \((\varphi_{n}f,(g_{1}+g_{2})\varphi_{n}f)\) 的类对（它们属于 \(m_{g_{1}}+m_{g_{2}}\) 的图像）所成的序列，收敛于 \((f,h)\) 在 \(\mathrm{L}^{2}(\mathrm{X},\mu)\) 中的类对。因此 \(m_{g_{1}}+m_{g_{2}}\) 的闭包确实等于 \(m_{g_{1}+g_{2}}\)。

\(m_{g_{1}} \circ m_{g_{2}} \subset m_{g_{1}g_{2}}\) 是初等的。设 \(g_{1}\) 有界，于是 \(m_{g_{1}}\) 是 \(\mathrm{L}^{2}(\mathrm{X},\mu)\) 上的连续线性映射。此时 \(m_{g_{2}} \circ m_{g_{1}}\) 的定义域是这样的函数 \(f \in \mathcal{L}^{2}(\mathrm{X},\mu)\) 的等价类所成的集合：这些函数使 \(g_{2}(g_{1}f) \in \mathcal{L}^{2}(\mathrm{X},\mu)\)，亦即 \(m_{g_{1}g_{2}}\) 的定义域。因此我们有 \(m_{g_{2}} \circ m_{g_{1}} = m_{g_{1}g_{2}}\)。

同样初等的是：\(\operatorname{dom}(m_{g_{1}} \circ m_{g_{2}}) = \operatorname{dom}(m_{g_{2}})\) 包含于 \(\operatorname{dom}(m_{g_{1}g_{2}})\)，且 \(m_{g_{1}g_{2}}\) 到这个空间的限制等于 \(m_{g_{1}} \circ m_{g_{2}}\)。

应当警惕，不要以为一般地有 \(m_{g_{1}} \circ m_{g_{2}} = m_{g_{1}g_{2}}\)（IV，第 347 页，习题 10）。尽管如此，我们有如下的部分结果：

命题 24. --- 设 g 是 X 上的 μ-可测函数。

\begin{enumerate}
\def\labelenumi{\alph{enumi})}
\item
  有 \(m_{\overline{g}} \circ m_g = m_{|g|^2}\)；
\item
  对一切整数 \(k, \ell \in \mathbf{N}\)，有 \(m_{g^k\overline{g}^\ell} = m_g^k m_{\overline{g}}^\ell\)。
\end{enumerate}

由引理 6 有 \(m_{\overline{g}} \circ m_g \subset m_{|g|^2}\)。反之，由不等式 \(|g| \leqslant 1 + |g|^2\) 推出 \(\text{dom}(m_{|g|^2}) = \text{dom}(m_{\overline{g}} \circ m_g)\)，由此得第一个断言。

设 \(k, \ell \in \mathbf{N}\)。我们有 \(m_g^k m_{\overline{g}}^\ell \subset m_{g^k \overline{g}^\ell}\)（同上所引）。\(m_g^k m_{\overline{g}}^\ell\) 的定义域是在 \(\mathrm{L}^2(\mathrm{X}, \mu)\) 中这样的函数 \(h \in \mathcal{L}^2(\mathrm{X}, \mu)\) 的类所成的集合：\(|g|^j h\) 属于 \(\mathcal{L}^2(\mathrm{X}, \mu)\)，对每个整数 \(j\)（满足 \(0 \leqslant j \leqslant k + \ell\)）。不等式

\[
\left| g ^ {j} h \right| \leqslant | h | + \left| g ^ {k + \ell} h \right|,
\]

对 \(0 \leqslant j \leqslant k + \ell\) 成立，由它可以看出 \(\mathrm{dom}(m_g^k m_{\overline{g}}^\ell)\) 等于 \(\mathrm{dom}(m_{g^k\overline{g}^\ell})\)，由此得断言 \(b\)）。

